\emph{Proof.}
Fix one of the two profiles in the statement and denote it by \(w_d\).  For the arm-normalized choice, positivity and finiteness follow from Assumption~\(\mathrm{ass{:}nondegenerate\text{-}components}\), positive costs, and Theorem~\(\mathrm{thm{:}budget\text{-}profile}\).  For the IKM-normalized choice, Definition~\(\mathrm{def{:}ikmw\text{-}normalized\text{-}design\text{-}value}\) gives
\[
 V_{g,d}^{\mathrm{IKM}}=(1-p_{g,1})^2V_{g,d}.
\]
Assumption~\(\mathrm{ass{:}source\text{-}arm\text{-}positivity}\) makes each factor \(1-p_{g,1}\) strictly positive, and Assumption~\(\mathrm{ass{:}nondegenerate\text{-}components}\) makes \(V_{g,d}\) finite and strictly positive.  Consequently
\[
 0<w_d<\infty
 \qquad\text{for every }d\in\mathcal D                 \tag{10}
\]
under either profile.  The finite nonempty family \(\mathcal D\) therefore has nonempty minimizer sets, and every ratio below is well-defined.

Put \(r_d=R_d^A(P)\) and
\[
 \mathcal O_P^A=\{(d,d')\in\mathcal D^2:r_d\le r_{d'}\}.
\]
This set is finite and contains all diagonal pairs.  For \(C\ge0\), all treatment-aware order implications hold exactly when
\[
 \frac{w_d}{w_{d'}}\le C
 \qquad\text{for every }(d,d')\in\mathcal O_P^A.          \tag{11}
\]
By finiteness, equation~\((11)\) is equivalent to
\[
 C\ge
 \max_{(d,d')\in\mathcal O_P^A}\frac{w_d}{w_{d'}},       \tag{12}
\]
which proves equation~\((1)\) and its least-constant assertion.

Choose \(d_{\mathrm{sel}}\) to minimize \(w_d\) over the treatment-aware prediction minimizers \(\mathcal D_{\mathrm{pred}}^A(P)\), and choose \(d_*\) to minimize \(w_d\) over all of \(\mathcal D\).  Then
\[
 r_{d_{\mathrm{sel}}}\le r_{d_*},
 \qquad
 1\le\frac{w_{d_{\mathrm{sel}}}}{w_{d_*}}
 \le \max_{(d,d')\in\mathcal O_P^A}\frac{w_d}{w_{d'}}.  \tag{13}
\]
For \(w_d=\mathcal V_d^*(P)\), the middle ratio in equation~\((13)\) is exactly \(\rho_{\mathrm{pred}}^A(P)\) by Definition~\(\mathrm{def{:}predictive\text{-}benchmark}\).  For \(w_d=\mathcal V_{d,\mathrm{IKM}}^*(P)\), it is exactly \(\rho_{\mathrm{pred}}^{A,\mathrm{IKM}}(P)\) by Definition~\(\mathrm{def{:}ikmw\text{-}normalized\text{-}design\text{-}value}\).  Thus each named ratio is finite at the fixed law and is bounded by equation~\((1)\).  Over a nonempty law class, any finite uniform constant bounds the corresponding exact ratio; conversely, the finite supremum of those ratios is itself a uniform constant.  This proves the stated uniform equivalence.

For the numerical envelope, define
\[
 s_d(x)=\frac{F_d(x)}{Q_{R,d}^A(x)}.
\]
Positivity of the prediction form on the chosen nonzero cone makes the denominator positive.  The definitions of \(\Lambda_-^A(P)\) and \(\Lambda_+^A(P)\) therefore imply
\[
 \Lambda_-^A(P)Q_{R,d}^A(x)
 \le F_d(x)
 \le\Lambda_+^A(P)Q_{R,d}^A(x)                          \tag{14}
\]
for every represented \(d\) and every cone direction \(x\).  Evaluating equation~\((14)\) at \(x_d(P)\) gives
\[
 \Lambda_-^A(P)R_d^A(P)
 \le\mathcal V_d^*(P)
 \le\Lambda_+^A(P)R_d^A(P).                              \tag{15}
\]
If \(R_d^A(P)\le R_{d'}^A(P)\), then equation~\((15)\) yields
\[
 \frac{\mathcal V_d^*(P)}{\mathcal V_{d'}^*(P)}
 \le
 \frac{\Lambda_+^A(P)R_d^A(P)}
      {\Lambda_-^A(P)R_{d'}^A(P)}
 \le\frac{\Lambda_+^A(P)}{\Lambda_-^A(P)}.              \tag{16}
\]
Taking the maximum over the treatment-aware order relation and using equation~\((13)\) proves equation~\((2)\).

Finally, Theorem~\(\mathrm{thm{:}sharp\text{-}predictive\text{-}operator\text{-}certificate}\) proves algebraically that profiling the two source-score quadratic forms produces
\[
 F_d(x)=\inf_{0<\alpha<1}
 \left\{
 \frac{c_{E,d}Q_{E,d}(x)}\alpha+
 \frac{c_{O,d}Q_{O,d}(x)}{1-\alpha}
 \right\},                                                \tag{17}
\]
which is not a quadratic form in general.  Its explicit parallelogram-law counterexample applies unchanged after replacing the prediction denominator by \(Q_{R,d}^A\).  Moreover, equality in equation~\((16)\) requires simultaneous attainment of the upper and lower envelope endpoints by law-realized directions and saturation of the treatment-aware risk order.  A largest generalized eigenvalue at fixed \(\alpha\) supplies only an upper full-space quadratic envelope and does not supply these conditions.  Hence the two-sided certificate can be conservative and the sharp treatment-aware selection constant is not generally a largest generalized eigenvalue. \(\square\)